\documentclass[11pt,a4paper]{article}
\usepackage[T1]{fontenc}
\usepackage[margin=25mm]{geometry}
\usepackage{amsmath,amssymb,amsthm,mathtools,microtype}
\usepackage[hidelinks]{hyperref}
\allowdisplaybreaks
\setlength{\emergencystretch}{3em}
\newcommand{\C}{\mathbb C}
\newcommand{\R}{\mathbb R}
\newcommand{\im}{\operatorname{im}}
\newcommand{\lc}{\operatorname{lc}}
\newcommand{\diag}{\operatorname{diag}}
\newcommand{\norm}[1]{\left\lVert#1\right\rVert}
\newtheorem{theorem}{Theorem}
\newtheorem{proposition}[theorem]{Proposition}
\title{Independent derivation of the original observation update,\protect\\
its quotient current, and the two low-cutoff receivers}
\author{Split-Zero research programme}
\date{19 September 2026}
\begin{document}\maketitle

\section{Original objects, conventions, and the finite minimum}

This derivation checks ORE1--23 of
\texttt{OBSERVATION\_RESPONSE\_EVOLUTION.tex} directly from the
original polynomial source. It also derives further exact receiving
identities for equations (5)--(7) of the supplied
\texttt{LOW\_CUTOFF\_ENERGIES\_CONDUCTOR\_AND\_MARKED\_DEFECT.md}.
All adjoints below are conjugate transposes in the stated fixed frames;
the inner product is antilinear in its first argument. No original
class, source mass, period map, relation, or conductor coordinate is
discarded. The calculation is finite and uses the original source
domain on which the polynomial norms below are defined.

Keep the actual simple-quartet source, including its original unit,
period and logarithm branch, and put
\begin{equation}\label{oid:objects}
\begin{gathered}
a\equiv1\pmod4,\qquad q=(a+1)^2,\qquad c=a/2,\qquad
0<\delta<1/2,\quad\gamma>2,\\
\chi(S)=\prod_{u,b=0}^a
\bigl[S-c-(2u-a)\delta-i(2b-a)\gamma\bigr],\qquad
E=\C[S]/(\chi),\\
dm_a(y)=w_h^{*a}(y)\,dy,\qquad
w_h(y)=\frac1{2\pi}\left|\frac{2\xi}{h}(1/2+iy)\right|^2,\\
\langle f,g\rangle_{m_a}
=\int_{\R}\overline{f(c+iy)}g(c+iy)\,dm_a(y).
\end{gathered}
\end{equation}
The frame of $E$ is the original increasing-power frame. The original
source gives a positive inner product on each finite polynomial space:
a nonzero polynomial is nonzero except at finitely many points, and
the original source density is nonzero on intervals. The moments used
in the original finite source norms are finite. These are the retained
source hypotheses, not a replacement Gamma norm.

Write $\mathcal P_N=\C[S]_{\leq N}$ and
$\mathcal R_N=\chi\mathcal P_{N-q}$, with $\mathcal P_j=0$ for
$j<0$. For $N\geq q-1$, let $\mathcal V_N$ be the orthogonal
complement of $\mathcal R_N$ in $\mathcal P_N$. Polynomial division
by the monic $\chi$ gives a surjection
$\pi_N:\mathcal P_N\to E$ whose kernel is precisely $\mathcal R_N$.
Consequently $\pi_N|_{\mathcal V_N}$ is bijective: its kernel is
zero and both spaces have dimension $q$. Denote its inverse by
$T_N:E\to\mathcal V_N$. Every lift of $x\in E$ is uniquely
$T_Nx+r$ with $r\in\mathcal R_N$, and
\[
 \norm{T_Nx+r}_{m_a}^2=\norm{T_Nx}_{m_a}^2+\norm r_{m_a}^2.
\]
Thus $T_Nx$ is the unique attained minimum over the complete relation
space and $G_N$ is exactly its positive Hermitian metric:
$x^*G_Ny=\langle T_Nx,T_Ny\rangle_{m_a}$.

Let $p_n(S)$ be the monic source orthogonal polynomial and
$\omega_n=\norm{p_n}_{m_a}^2>0$, and put $b_n=[p_n]\in E$.
The polynomials $p_0,\ldots,p_N$ form a basis, because their distinct
degrees and leading coefficients give a triangular change of basis
with diagonal one. Define
\[
 C_N=[b_0,\ldots,b_N],\qquad
 \Omega_N=\diag(\omega_0,\ldots,\omega_N),\qquad
 M_N=C_N\Omega_N^{-1}C_N^*.
\]
The matrix $C_N$ is surjective, so $M_N$ is positive definite:
$z^*M_Nz=\sum_{j=0}^N|b_j^*z|^2/\omega_j$ vanishes only when
$C_N^*z=0$, hence only when $z=0$. For $x\in E$ the coefficient
vector
\[
 t_N(x)=\Omega_N^{-1}C_N^*M_N^{-1}x
\]
satisfies $C_Nt_N(x)=x$. If $C_Nu=0$, then
$u^*\Omega_Nt_N(x)=u^*C_N^*M_N^{-1}x=0$.
Expanding the positive square of $t_N(x)+u$ proves its minimality
and yields the explicit, original-source formula
\begin{equation}\label{oid:minimum}
 G_N=M_N^{-1},\qquad
 T_Nx=\sum_{j=0}^N p_j\,
 \frac{b_j^*G_Nx}{\omega_j},\qquad
 G_N^{-1}=\sum_{j=0}^N\frac{b_jb_j^*}{\omega_j}.
\end{equation}
This derives the dual norm and the attained minimum simultaneously.

\section{The original source increment and the full boundary action}

At one step write, only to display the original factors compactly,
\[
 G=G_N,\quad G'=G_{N+1},\quad b=b_{N+1},\quad
 \omega=\omega_{N+1},\quad \beta=b^*Gb,\quad D=\omega+\beta.
\]
The new summand in \eqref{oid:minimum} is exactly $bb^*/\omega$.
Direct multiplication, using $b^*Gb=\beta$, proves
\begin{equation}\label{oid:rankone}
 (G')^{-1}=G^{-1}+\frac{bb^*}{\omega},\qquad
 G'=G-\frac{Gbb^*G}{D},\qquad
 G'b=rGb,\quad b^*G'b=r\beta,
 \qquad r=\frac{\omega}{D}.
\end{equation}
Indeed the coefficient of $Gbb^*$ in the product of the proposed
$G'$ with $G^{-1}+bb^*/\omega$ is
$\omega^{-1}-D^{-1}-\beta/(D\omega)=0$.
The resulting inverse identity is the rank-one update recorded by
Hager~\cite[p.~221, equation~(2)]{Hager}; its source column and its
positive denominator have been derived here from the actual minimum.
For a column $x$ and a row $y^*$, multilinearity of determinant
implies $\det(I+xy^*)=1+y^*x$: all terms using two columns from
$xy^*$ vanish, and the one-column terms sum to its trace.
Applying this identity to \eqref{oid:rankone} gives
\begin{equation}\label{oid:r}
 \frac{\det G'}{\det G}=1-\frac\beta D
 =\frac\omega D=r.
\end{equation}

We next prove the boundary formula needed for the currents, without
importing a conjugation convention from a response formula. Let $A:E\to E$
be multiplication by $S$. For $x\in E$, the coefficient of $S^N$ in
$T_Nx$ is
\begin{equation}\label{oid:leading}
 \ell_Nx=\frac{b_N^*G_Nx}{\omega_N}.
\end{equation}
This follows either from \eqref{oid:minimum}, since $p_N$ is monic,
or by pairing $p_N$ with $T_Nx$: the difference $p_N-T_Nb_N$
is an old relation and is orthogonal to $T_Nx$.

Let $\Pi_N$ be orthogonal projection onto $\mathcal V_N$, acting
also on $\mathcal P_{N+1}$, and let
$J_N=\Pi_N M_y|_{\mathcal V_N}$. It is Hermitian, since
multiplication by the real variable $y$ is Hermitian in the original
integral. Set $\widehat b_j=T_Nb_j$. For $f=T_Nx$,
$Sf-(\ell_Nx)p_{N+1}$ has degree at most $N$ and class
$Ax-(\ell_Nx)b_{N+1}$. Since $p_{N+1}\perp\mathcal P_N$,
projecting this polynomial onto $\mathcal V_N$ gives
\begin{equation}\label{oid:action}
 \begin{gathered}
 T_NAT_N^{-1}=cI+iJ_N+
 \frac{\widehat b_{N+1}\widehat b_N^*}{\omega_N},\\
 D_N:=A^*G_N+G_NA-aG_N
 =\frac{G_Nb_{N+1}b_N^*G_N+G_Nb_Nb_{N+1}^*G_N}{\omega_N}.
 \end{gathered}
\end{equation}
Here the star in the source-space rank-one operator denotes its
actual source inner product; the matrix identity on $E$ uses the
original power frame. The sign is positive because the class
$(\ell_Nx)b_{N+1}$ is added back after the degree reduction.

Retain the full original eigenclass
\begin{equation}\label{oid:eigen}
 \lambda=c+d+it,\quad d=a\delta>0,\quad t=a\gamma>0,
 \qquad v=v_\lambda=\frac{\chi(S)}{(S-\lambda)\chi'(\lambda)},
 \qquad Av=\lambda v,
\end{equation}
including all its original low, conductor and native coordinates.
Put $f=T_Nv$, $E_N=v^*G_Nv$, and
$\mu_N=E_N^{-1}\int y|f(c+iy)|^2dm_a(y)\in\R$.
The polynomial $(S-\lambda)f$ is a relation. If
$r_N(S)=\chi(S)S^{N+1-q}$, its difference from
$(\ell_Nv)r_N$ is an old relation. Moreover
$p_{N+1}-r_N\in\mathcal P_N$ represents $b_{N+1}$, so
\[
 \langle f,r_N\rangle=-\langle f,p_{N+1}-r_N\rangle
 =-v^*G_Nb_{N+1}.
\]
Consequently
\[
 E_N[-d+i(\mu_N-t)]
 =\langle f,(S-\lambda)f\rangle
 =-\frac{\overline{b_{N+1}^*G_Nv}\,b_N^*G_Nv}{\omega_N}.
\]
Taking the conjugate, with the displayed minus sign retained, proves
\begin{equation}\label{oid:current}
 \mathfrak c_N
 :=\frac{\overline{b_N^*G_Nv}\,b_{N+1}^*G_Nv}
 {\omega_NE_N}
 =d+i(\mu_N-t),\qquad \Re\mathfrak c_N=d>0.
\end{equation}
Both boundary inner products are therefore nonzero. In particular
$b\ne0$, $\beta>0$, and $0<r<1$ in \eqref{oid:rankone}.

The original measure is even and
$\chi(a-S)=(-1)^q\chi(S)$, directly from the paired full root grid.
The isometry $f(S)\mapsto f(a-S)$ preserves both the relation space
and its orthogonal complement. Uniqueness of the monic orthogonal
polynomial gives $p_j(a-S)=(-1)^jp_j(S)$, so the attained columns
$\widehat b_N,\widehat b_{N+1}$ have opposite parities. Hence
\begin{equation}\label{oid:parity}
 b_N^*G_Nb_{N+1}=0.
\end{equation}
The squared norm of the projection of $v$ onto these two nonzero
orthogonal columns is
\[
 \frac{|b_N^*G_Nv|^2}{b_N^*G_Nb_N}
 +\frac{|b_{N+1}^*G_Nv|^2}{b_{N+1}^*G_Nb_{N+1}}\le E_N.
\]
Thus the original two boundary energies satisfy
\begin{equation}\label{oid:energies}
 p_{1,N}:=\frac{|b_N^*G_Nv|^2}{E_Nb_N^*G_Nb_N}>0,\qquad
 p_{2,N}:=\frac{|b_{N+1}^*G_Nv|^2}{E_Nb_{N+1}^*G_Nb_{N+1}}>0,
 \qquad p_{1,N}+p_{2,N}\le1.
\end{equation}
In particular each energy is strictly below one on this original
eigenclass domain.

\section{The original observation minimum and its exact update}

Fix the full original observation $\Lambda:E\to B=\im\Lambda$
in a fixed frame of $B$, retaining its period map. Let $K=\ker\Lambda$
and let $I:\C^{\dim K}\to E$ be any fixed injective frame of $K$.
At the present cutoff define
\begin{equation}\label{oid:observation}
 H=I^*GI,\qquad P=IH^{-1}I^*G,\qquad
 Q=(\Lambda G^{-1}\Lambda^*)^{-1},\qquad
 L=G^{-1}\Lambda^*Q:B\to E.
\end{equation}
Surjectivity of $\Lambda$ makes $\Lambda G^{-1}\Lambda^*$
positive definite. The identity $P^2=P$, its image $K$, and
$P^*G=GP$ follow by multiplication. Also $\Lambda L=I_B$ and
$I^*GL=I^*\Lambda^*Q=0$. Every lift of $y\in B$ is uniquely
$Ly+k$ with $k\in K$ and its squared norm is
$y^*Qy+k^*Gk$. Thus $L$ is the actual unique minimum lift and
\begin{equation}\label{oid:quotient}
 L\Lambda=I_E-P,\qquad
 G(I_E-P)=\Lambda^*Q\Lambda.
\end{equation}
For $K=0$ the corresponding inverse in $H$ is empty and $P=0$;
for $B=0$ the corresponding inverse in $Q$ is empty. Empty
determinants equal one throughout.

Keep
\[
 \alpha=\frac{b^*GPb}{\beta}\in[0,1],\qquad
 \xi=1-r=\frac\beta D,\qquad
 d_K=\frac{\det H'}{\det H},\qquad u=\Lambda b.
\]
Restriction of \eqref{oid:rankone} to $I$ and the proved determinant
identity give
\[
 H'=H-\frac{zz^*}{D},\qquad z=I^*Gb,\qquad
 z^*H^{-1}z=\beta\alpha,
\]
and hence
\begin{equation}\label{oid:allocations}
 d_K=1-\xi\alpha,\qquad r\le d_K\le1,\qquad
 D^B:=\omega+\beta(1-\alpha)=Dd_K>0.
\end{equation}
The quotient inverse updates by the actual observed column:
\[
 (Q')^{-1}=Q^{-1}+\frac{uu^*}{\omega},\qquad
 u^*Qu=b^*G(I-P)b=\beta(1-\alpha).
\]
Multiplication as in \eqref{oid:rankone} now proves
\begin{equation}\label{oid:Qupdate}
 Q'=Q-\frac{Quu^*Q}{D^B},\qquad
 r^B:=\frac{\det Q'}{\det Q}
 =\frac\omega{D^B}=\frac r{d_K}.
\end{equation}
In particular the exact nonnegative logarithmic allocations are
$-\log d_K$ and $-\log r+\log d_K$, with sum $-\log r$.

There is also an exact update of the original minimum section and
of the projector itself. Multiplying the two matrices in
$L'=(G^{-1}+bb^*/\omega)\Lambda^*Q'$ and using
$\omega+u^*Qu=D^B$ gives
\begin{equation}\label{oid:projector}
 L'=L+\frac{Pb\,u^*Q}{D^B},\qquad
 P'=P-\frac{Pb\,b^*G(I-P)}{D^B}.
\end{equation}
The second identity follows by multiplying the first by $\Lambda$
and applying \eqref{oid:quotient}. Its correction is an actual map
into the original kernel; it does not replace that kernel by a rank.
It is defined even when $u=0$ or $\Lambda v=0$.

\section{Response transport, conjugations, and energy transport}

Return cutoff indices to every quantity and retain the original rows
\begin{equation}\label{oid:UV}
 U_N=\frac{b_N^*G_NP_Nv}{b_N^*G_Nv},\qquad
 V_N=\frac{b_{N+1}^*G_NP_Nv}{b_{N+1}^*G_Nv},\qquad
 \theta_N=\frac{v^*G_NP_Nv}{E_N}.
\end{equation}
They have their original order. Write
$s_N=b_{N+1}^*G_Nv$ and
$s_N^B=b_{N+1}^*G_N(I-P_N)v=(1-V_N)s_N$.
From \eqref{oid:rankone}, the next full numerator for this same
$b_{N+1}$ is $r_{N+1}s_N$. From \eqref{oid:Qupdate}, its next
quotient numerator is
\begin{equation}\label{oid:transportnumerator}
 b_{N+1}^*G_{N+1}(I-P_{N+1})v
 =u^*Q_{N+1}\Lambda v
 =\frac{\omega_{N+1}}{D_N^B}s_N^B
 =\frac{r_{N+1}}{d_{K,N}}s_N^B.
\end{equation}
Dividing by the proved nonzero full numerator yields
\begin{equation}\label{oid:response}
 \boxed{U_{N+1}=1-\frac{1-V_N}{d_{K,N}}.}
\end{equation}
No conjugation occurs in this linear transport. Conjugation enters
only when the earlier row is used in a Hermitian current.

For $x_{K,N}=P_Nv$ and $x_{B,N}=(I-P_N)v$, define the actual
centered pairings
\[
 \Phi_{K,N}=x_{K,N}^*D_Nx_{K,N},\qquad
 \Phi_{B,N}=x_{B,N}^*D_Nx_{B,N},\qquad
 \Phi_{\times,N}=2\Re(x_{K,N}^*D_Nx_{B,N}).
\]
Equation \eqref{oid:action} gives, with every orientation explicit,
\begin{equation}\label{oid:products}
\begin{gathered}
 P_{K,N}^{\rm resp}=\overline U_NV_N,\qquad
 P_{B,N}^{\rm resp}=(1-\overline U_N)(1-V_N),\\
 P_{\times,N}^{\rm resp}=\overline U_N+V_N-2\overline U_NV_N,
 \qquad \sum_XP_{X,N}^{\rm resp}=1,\\
 \Phi_{X,N}=2E_N\Re(\mathfrak c_NP_{X,N}^{\rm resp}),\qquad
 \sum_X\Phi_{X,N}=2dE_N.
\end{gathered}
\end{equation}
For example its quotient term is
$2\Re(\overline{b_N^*G_Nx_{B,N}}\,
b_{N+1}^*G_Nx_{B,N})/\omega_N$. Its two cross terms are
$2\Re(\overline{b_N^*G_Nx_{K,N}}\,b_{N+1}^*G_Nx_{B,N}
+\overline{b_N^*G_Nx_{B,N}}\,b_{N+1}^*G_Nx_{K,N})/\omega_N$.
These expansions prove the three products, including their conjugations.
For $N\ge q$, put $R_N=1-V_N$. Then \eqref{oid:response} gives
\begin{equation}\label{oid:consecutiveproducts}
 P_{B,N}^{\rm resp}
 =\frac{\overline{R_{N-1}}R_N}{d_{K,N-1}},\qquad
 P_{K,N}^{\rm resp}
 =\left(1-\frac{\overline{R_{N-1}}}{d_{K,N-1}}\right)(1-R_N).
\end{equation}
At $N=q-1$ the original $U_{q-1}$ is retained; there is no earlier
surjective quotient source of dimension $q$ to insert.

Set $\eta_N=(1-r_{N+1})p_{2,N}$. Substitution into
\eqref{oid:rankone}, followed by \eqref{oid:energies}, proves
\begin{equation}\label{oid:Eloss}
 E_{N+1}=E_N-\frac{|s_N|^2}{\omega_{N+1}+\beta_N},\qquad
 \frac{E_{N+1}}{E_N}=1-\eta_N,
 \qquad 0<\eta_N<1-r_{N+1}<1,
 \qquad r_{N+1}<\frac{E_{N+1}}{E_N}<1.
\end{equation}
The strict comparison with $r_{N+1}$ strengthens the non-strict
Cauchy--Schwarz bound in ORE13; it follows from the other nonzero,
orthogonal boundary column, not from an assumed asymptotic profile.
For the observed energy
$e_{B,N}=E_N(1-\theta_N)=(\Lambda v)^*Q_N\Lambda v$,
\eqref{oid:Qupdate} yields
\begin{equation}\label{oid:observedloss}
 \Delta_N^B:=e_{B,N}-e_{B,N+1}
 =\frac{|s_N^B|^2}{D_N^B}\ge0,
 \qquad
 (1-\eta_N)(1-\theta_{N+1})
 =1-\theta_N-\frac{\eta_N|1-V_N|^2}{d_{K,N}}.
\end{equation}
Rearranging the second equation, without dividing by $1-\theta_N$,
gives
\begin{equation}\label{oid:theta}
 \theta_{N+1}-\theta_N
 =\frac{\eta_N}{1-\eta_N}
 \left[\theta_N-1+\frac{|1-V_N|^2}{d_{K,N}}\right],\qquad
 0\le\frac{\eta_N|1-V_N|^2}{d_{K,N}}\le1-\theta_N.
\end{equation}
If $\Lambda v\ne0$, positive definiteness of $Q_{N+1}$ makes the
last upper inequality strict. If $\Lambda v=0$, $v\in K$ and
$P_Nv=v$, so $\theta_N=1$, $U_N=V_N=1$, and every observed
energy and observed drop is zero. This evaluates that stratum without
an undefined logarithm or phase.

Finally \eqref{oid:rankone} gives, for the same column $b_{N+1}$,
\begin{equation}\label{oid:p1transport}
 p_{1,N+1}=
 \frac{r_{N+1}p_{2,N}}
 {1-(1-r_{N+1})p_{2,N}}.
\end{equation}
Both numerator and denominator use the retained class and original
source metric at their displayed cutoffs.

\section{The quotient-current receiver and its finite bound}

For $N\ge q$, \eqref{oid:transportnumerator} says that the actual
first quotient numerator at cutoff $N$ is
\[
 b_N^*G_N(I-P_N)v
 =\frac{\omega_N}{D_{N-1}^B}s_{N-1}^B.
\]
Insert this identity, with its real positive coefficient, into the
quotient expansion following \eqref{oid:products}. The factor
$\omega_N$ cancels its matching boundary denominator, proving
\begin{equation}\label{oid:quotientcurrent}
 \boxed{\Phi_{B,N}
 =\frac{2\Re(\overline{s_{N-1}^B}s_N^B)}{D_{N-1}^B}.}
\end{equation}
This is precisely ORE21. The conjugate belongs to $s_{N-1}^B$;
neither replacing $\overline{s_{N-1}^B}s_N^B$ by
$s_{N-1}^Bs_N^B$ nor replacing $D_{N-1}^B$ by $D_N^B$
preserves the displayed derivation. It agrees exactly with
\eqref{oid:consecutiveproducts}, because
$s_N^B=(1-V_N)s_N$ and
$b_N^*G_Nv=r_Ns_{N-1}$.

When both drops are positive, set $z_N=s_N^B/|s_N^B|$ and
$c_N^B=\sqrt{D_N^B/D_{N-1}^B}>0$. Equation
\eqref{oid:observedloss} then proves
\begin{equation}\label{oid:phase}
 \Phi_{B,N}=2c_N^B\sqrt{\Delta_{N-1}^B\Delta_N^B},
 \Re(\overline z_{N-1}z_N),\qquad
 |\Phi_{B,N}|\le2c_N^B\sqrt{\Delta_{N-1}^B\Delta_N^B}.
\end{equation}
If either drop is zero its numerator is zero and
\eqref{oid:quotientcurrent} gives $\Phi_{B,N}=0$; a phase at zero
is unnecessary. Since $D_N^B=\omega_{N+1}/r_{N+1}^B$,
\begin{equation}\label{oid:cB}
 (c_N^B)^2=\frac{\omega_{N+1}}{\omega_N},
 \frac{r_N^B}{r_{N+1}^B},\qquad
 \sum_{N=L}^H\Delta_N^B=e_{B,L}-e_{B,H+1}.
\end{equation}
For a finite interval $q\le L\le H$, applying
$2\sqrt{xy}\le x+y$ to each term of \eqref{oid:phase} gives
\begin{align}\label{oid:finitebound}
 \sum_{N=L}^H|\Phi_{B,N}|
 &\le \max_{L\le N\le H}c_N^B
 \left(\Delta_{L-1}^B+2\sum_{N=L}^{H-1}\Delta_N^B+\Delta_H^B\right)
 \notag\\
 &\le 2\max_{L\le N\le H}c_N^B,e_{B,L-1}.
\end{align}
The second inequality follows because all drops are nonnegative and
their sum from $L-1$ to $H$ is $e_{B,L-1}-e_{B,H+1}\le e_{B,L-1}$.
This proves ORE22--23 and their endpoint coefficients exactly. It
does not supply a uniform upper bound for the displayed $c_N^B$.

\section{A further exact relation: the next kernel overlap determines
the present quotient determinant}

Define the next first-boundary kernel overlap in its original metric,
\begin{equation}\label{oid:pi}
 \pi_{1,K,N+1}
 =\frac{b_{N+1}^*G_{N+1}P_{N+1}b_{N+1}}
 {b_{N+1}^*G_{N+1}b_{N+1}}\in[0,1].
\end{equation}
In the one-step notation, $z=I^*Gb$ and
$H'=H-zz^*/D$. Multiplication proves
\[
 (H')^{-1}=H^{-1}
 +\frac{H^{-1}zz^*H^{-1}}{D^B},\qquad
 z^*(H')^{-1}z
 =\beta\alpha+\frac{(\beta\alpha)^2}{D^B}
 =\frac{\beta\alpha}{d_K}.
\]
As $I^*G'b=rI^*Gb$, the numerator in \eqref{oid:pi} is
$r^2\beta\alpha/d_K$, while its denominator is $r\beta$.
Therefore, with $\pi=\pi_{1,K,N+1}$,
\begin{equation}\label{oid:newallocation}
 \boxed{\pi=\frac{r\alpha}{d_K},\qquad
 \alpha=\frac{\pi}{r+(1-r)\pi},\qquad
 d_K=\frac r{r+(1-r)\pi},\qquad
 r^B=r+(1-r)\pi.}
\end{equation}
All denominators are positive since $r>0$. For $K=0$ these give
$\pi=\alpha=0,d_K=1,r^B=r$; for $K=E$ they give
$\pi=\alpha=1,d_K=r,r^B=1$. Thus the endpoint conventions are
included in the same exact identity.

This relation is additional to ORE1--23. It makes a bound on the
actual next first-boundary kernel overlap into a bound on the same
step's quotient determinant. In particular the source's conductor
bound for $\pi_{1,K,q}$ has this exact receiving map at the first
source increment. It does not require a choice of kernel basis or a
replacement by the dimension of that kernel.

\section{The complete coefficient chart and both original windows}

For completeness the metric update above retains the full arithmetic
coordinate map. Write $\chi(S)=\sum_{\ell=0}^qc_\ell S^\ell$,
$c_q=1$, and define the residue-dual polynomials
\[
 \mathfrak p_n(S)=\sum_{\ell=n+1}^qc_\ell S^{\ell-n-1},
 \qquad 0\le n<q,
 \qquad F e_n=\mathfrak p_n.
\]
Their distinct monic degrees $q-1-n$ prove that $F:\C^q\to E$
is invertible. With $\mathbb V_{n,\alpha}=\beta_\alpha^n$ and
$\mathcal L(\nu)=\sum_\alpha\nu_\alpha\chi(S)/(S-\beta_\alpha)$,
the residue pairing gives $F=\mathcal L\mathbb V^{-1}$. To see the
pairing directly, the sum of finite residues of
$\mathfrak p_n(S)S^j/\chi(S)$ is its $S^{-1}$ coefficient at infinity.
It is zero for $j<n$, one for $j=n$, and zero for $j>n$, the latter
from
$S^{n+1}\mathfrak p_n=\chi-\sum_{\ell=0}^nc_\ell S^\ell$.
This proves the claimed inverse moment map with its ordinary
Vandermonde transpose convention.

Use $v_A=\operatorname{ord}_0 E_A$ for the conductor vanishing
order, to avoid confusing it with the retained class $v$.
Put $h_A=q-v_A$ and $d_A=(a-7)^2$, on the original conductor
domain. Let $P_{v_A}$ insert the initial $v_A$ zero coordinates.
The actual conductor map $C_a^{\mathsf T}$ satisfies
\[
 \mathbb V C_a^{\mathsf T}=P_{v_A}J,
 \qquad
 J_{i,\alpha'}=\sum_{r=0}^i\binom{v_A+i}{r}
 \mu_{v_A+i-r}(\beta'_{\alpha'})^r.
\]
This is the ordinary product rule for
$E_A(z)\sum_{\alpha'}\eta_{\alpha'}e^{\beta'_{\alpha'}z}$.
The map $C_a^{\mathsf T}$ is injective: if this entire product is
zero, the nonzero analytic factor $E_A$ is nonzero on an open disc,
so the exponential polynomial vanishes identically; its first $d_A$
derivatives give an invertible Vandermonde and force $\eta=0$.
Thus $J$ has rank $d_A$. Retain an actual selected minor
$I_{\rm sel}$ with $J_{I_{\rm sel}}$ invertible and complement
$J(I_{\rm sel})$. Put
\[
 H_{I_{\rm sel}}=J_{I_{\rm sel}}^{-1}E_{I_{\rm sel}}^{\mathsf T},
 \qquad
 Q_{I_{\rm sel}}=E_{J(I_{\rm sel})}^{\mathsf T}
 -J_{J(I_{\rm sel})}H_{I_{\rm sel}}.
\]
Checking the selected and complementary rows separately proves
$JH_{I_{\rm sel}}+E_{J(I_{\rm sel})}Q_{I_{\rm sel}}=I$.
It follows that the full chart
\[
 \mathcal F=[F_0,F_C,L_{I_{\rm sel}}],\qquad
 F_0=FE_{\{0,\ldots,v_A-1\}},\quad
 F_C=FP_{v_A}J=\mathcal LC_a^{\mathsf T},\quad
 L_{I_{\rm sel}}=FP_{v_A}E_{J(I_{\rm sel})}
\]
is invertible. Its inverse retains all three coordinates, including
$H_{I_{\rm sel}}P_{v_A}^{\mathsf T}F^{-1}x$ in the conductor block.
For the consecutive chart $0$, application of its $H_0,Q_0$ to the
selected coefficient columns proves the complete transition
\[
 L_{I_{\rm sel}}=F_CD_{0I_{\rm sel}}+L_0M_{0I_{\rm sel}},\quad
 D_{0I_{\rm sel}}=H_0E_{J(I_{\rm sel})},\quad
 M_{0I_{\rm sel}}=Q_0E_{J(I_{\rm sel})}.
\]
No change of the observation frame $I$ is implicit in this native
section notation. Set $Y=\mathcal F^{-1}I$. Then
\begin{equation}\label{oid:fullchart}
 \Lambda\mathcal FY=0,\qquad
 H_N=Y^*\mathcal F^*G_N\mathcal FY.
\end{equation}
If $Y=(Y_0,Y_C,Y_I)^{\mathsf T}$, the latter is explicitly the
sum of all nine terms
$\sum_{r,s\in\{0,C,I\}}Y_r^*F_r^*G_NF_sY_s$.
For example the native diagonal block under its transition is
\[
 D^*F_C^*G_NF_CD+D^*F_C^*G_NL_0M
 +M^*L_0^*G_NF_CD+M^*L_0^*G_NL_0M.
\]
Thus ORE19 retains every cross metric and every conductor term.
Any analytic-native formula using an ordinary transpose remains in
its own coefficient maps; this Hermitian pullback does not change
that transpose to an adjoint.

For every positive scalar sequence $f_N$, cancellation of consecutive
ratios proves
\begin{equation}\label{oid:windows}
 \log\frac{f_{2q-1}f_{2q}}{f_{q-1}f_q}
 =\sum_{N=q-1}^{2q-1}w_N\log\frac{f_{N+1}}{f_N},\quad
 w_{q-1}=w_{2q-1}=1,\quad w_N=2\quad(q\le N\le2q-2).
\end{equation}
Indeed the interior exponent of each $f_j$ is $w_{j-1}-w_j$;
the only surviving exponents are $-1,-1,+1,+1$ at
$q-1,q,2q-1,2q$. Use this with $\det G_N,\det H_N,\det Q_N,E_N$
and the increments $r_{N+1},d_{K,N},r_{N+1}/d_{K,N},1-\eta_N$.
It also applies to $e_{B,N}$ when $\Lambda v\ne0$. When
$\Lambda v=0$, that sequence is identically zero and its receiving
equation is \eqref{oid:observedloss}, not a logarithm. This proves
ORE18 with its original two windows and both endpoints.

\section{Exact low-cutoff evaluation, with all complex phases}

The supplied low-cutoff note uses $k$ for the same tensor degree
called $a$ above. In this section set $k=a$ and preserve the note's
$q=(k+1)^2$, $c=k/2$, $d=k\delta$, $t=k\gamma$ and
\[
 z=\frac{\lambda-c}{i}=t-id,\qquad
 Q(y)=i^{-q}\chi(c+iy),\qquad b_\lambda=\chi'(\lambda)^{-1}.
\]
The full physical polynomial $Q$ is real and has parity $(-1)^q$.
Retain the original source integrals and their original masses:
\begin{equation}\label{oid:moments}
\begin{gathered}
 \nu_0=\int Q(y)^2dm_k(y),\qquad
 \nu_1=\int y^2Q(y)^2dm_k(y),\qquad
 d\mu_Q(y)=\nu_0^{-1}Q(y)^2dm_k(y),\\
 a_z=\int\frac{d\mu_Q(y)}{y-z},\quad
 m_z=\int\frac{d\mu_Q(y)}{|y-z|^2},\quad
 s_2=\frac{\nu_1}{\nu_0},\qquad
 d_0=\frac{\omega_q}{\nu_0},\quad
 d_1=\frac{\omega_{q+1}}{\nu_1},\\
 M=m_z-|a_z|^2,\qquad g=1+za_z.
\end{gathered}
\end{equation}
The probability $\mu_Q$ is only the note's auxiliary device for
ratios: every norm below retains the multiplying $\nu_0$ or $\nu_1$.
The symbols $d_0,d_1$ here are the two monic norm ratios of its
equation (4); they are not the unrelated tail coefficient $d_0(h)$
mentioned later in that note.

At $N=q-1$ no relation occurs in the minimum fibre, and its actual
polynomial representative is
\[
 f_{q-1}(c+iy)=b_\lambda i^{q-1}\frac{Q(y)}{y-z},\qquad
 E_{q-1}=|b_\lambda|^2\nu_0m_z.
\]
Its leading coefficient in $S$ is $b_\lambda$, so
$F_{13,q-1}=b_\lambda\omega_{q-1}$ and
$F_{11,q-1}=\omega_{q-1}$. The representative of $b_q$ is
$p_q-\chi$. Orthogonality gives
\[
 \norm{p_q-\chi}^2=\nu_0-\omega_q,
 \qquad
 \langle p_q-\chi,f_{q-1}\rangle=-\langle\chi,f_{q-1}\rangle.
\]
Since $\overline{i^q}i^{q-1}=i^{-1}=-i$, the latter numerator is
exactly $i b_\lambda\nu_0a_z$, including its phase. Thus
\begin{equation}\label{oid:lowfirst}
 \begin{gathered}
 F_{13,q-1}=b_\lambda\omega_{q-1},\qquad
 F_{23,q-1}=i b_\lambda\nu_0a_z,\qquad
 F_{22,q-1}=\nu_0-\omega_q,\\
 p_{1,q-1}=\frac{\omega_{q-1}}{\nu_0m_z},\qquad
 p_{2,q-1}=\frac{|a_z|^2}{(1-d_0)m_z},\qquad
 r_q=d_0.
 \end{gathered}
\end{equation}
The first two energy formulas rederive the note's equation (5).

At $N=q$ the relation space is $\C\chi$.
Its projection coefficient is
$\langle\chi,f_{q-1}\rangle/\nu_0=-i b_\lambda a_z$.
Subtracting this entire original relation gives
\begin{equation}\label{oid:lowsecondlift}
 f_q(c+iy)=b_\lambda i^{q-1}Q(y)
 \left(\frac1{y-z}-a_z\right),\qquad
 E_q=|b_\lambda|^2\nu_0M.
\end{equation}
This is the note's equation (6), obtained by the actual orthogonal
projection. Its first boundary numerator is, by the already proved
update, $F_{13,q}=d_0F_{23,q-1}=i b_\lambda\omega_q a_z$.

The next monic relation is $R(S)=\chi(S)(S-c)$.
It is orthogonal to $\chi$, because $\mu_Q$ is even, and its norm
is $\nu_1$. Both $p_{q+1}$ and $R$ have parity $(-1)^{q+1}$
about $c$; therefore $p_{q+1}-R$ is orthogonal to $\chi$ and is
the attained representative of $b_{q+1}$ at cutoff $q$.
Its norm is $\nu_1-\omega_{q+1}$. Moreover
\[
 \int\frac{y}{y-z}\,d\mu_Q(y)=1+za_z=g,\qquad
 \overline{i^{q+1}}i^{q-1}=i^{-2}=-1.
\]
It follows that
\begin{equation}\label{oid:lowsecond}
\begin{gathered}
 F_{13,q}=i b_\lambda\omega_q a_z,\qquad
 F_{23,q}=b_\lambda\nu_0g,\qquad
 F_{22,q}=\nu_1-\omega_{q+1},\qquad r_{q+1}=d_1,\\
 p_{1,q}=\frac{d_0|a_z|^2}{(1-d_0)M},\qquad
 p_{2,q}=\frac{|g|^2}{s_2(1-d_1)M}.
\end{gathered}
\end{equation}
This proves the note's equation (7) and its original phases. The
nonzero boundary numerators proved in \eqref{oid:current} give
$0<d_0,d_1<1$.

The same exact source update now determines further finite quantities:
\begin{equation}\label{oid:loweta}
 \boxed{\eta_{q-1}=\frac{|a_z|^2}{m_z},\qquad
 \eta_q=\frac{|g|^2}{s_2M},\qquad
 E_{q+1}=|b_\lambda|^2\nu_0
 \left(M-\frac{|g|^2}{s_2}\right).}
\end{equation}
For direct verification, at cutoff $q+1$ the entire relation space
is the orthogonal span of $\chi$ and $R$, and
\[
 f_{q+1}(c+iy)=b_\lambda i^{q-1}Q(y)
 \left(\frac1{y-z}-a_z-\frac{g}{s_2}y\right).
\]
Its original norm is the last expression in \eqref{oid:loweta}.
It is strictly positive because $1/(y-z)$ cannot equal an affine
function on a set containing an interval: multiplying such an
identity by $y-z$ would give a polynomial identity whose quadratic
and linear coefficients force both affine coefficients to be zero,
contradicting its constant coefficient one. This also proves the
strict inequalities $M>0$ and $M-|g|^2/s_2>0$ by their respective
orthogonal residuals. There is no missing new minimum in
\eqref{oid:loweta}.

The original full currents consequently have the explicit phases
\begin{equation}\label{oid:lowcurrent}
 \mathfrak c_{q-1}=\frac{i a_z}{m_z},\qquad
 \mathfrak c_q=\frac{-i\overline{a_z}g}{M},\qquad
 \boxed{\mathfrak c_q=
 \frac{\overline{\mathfrak c_{q-1}}
 -\eta_{q-1}(d+it)}{1-\eta_{q-1}}.}
\end{equation}
Indeed $\Im a_z=-dm_z$ by the integrand
$(y-z)^{-1}=(y-t-id)/|y-z|^2$, so the first current has real
part $d$. Also
$\Im(\overline{a_z}g)=-\Im a_z+|a_z|^2\Im z=dM$,
which proves the second current has real part $d$.
For its transport formula use
$a_z=-im_z\mathfrak c_{q-1}$,
$|a_z|^2/m_z=\eta_{q-1}$, and $iz=d+it$.
Equivalently, the original source means satisfy the exact relation
\begin{equation}\label{oid:lowmean}
 \mu_q=2t-\frac{\mu_{q-1}}{1-\eta_{q-1}}.
\end{equation}
It follows by comparing imaginary parts in \eqref{oid:lowcurrent}
and \eqref{oid:current}. These equations contain the original
physical center and imaginary coordinate; no phase is normalized away.

For the original observation, write
$\kappa_0=d_{K,q-1}$ and $\kappa_1=d_{K,q}$ only in the following
low-cutoff formulas. As $D_{q-1}=\nu_0$ and $D_q=\nu_1$,
\begin{equation}\label{oid:lowobservation}
\begin{gathered}
 D_{q-1}^B=\nu_0\kappa_0,\qquad
 D_q^B=\nu_1\kappa_1,\qquad
 (c_q^B)^2=s_2\frac{\kappa_1}{\kappa_0},\\
 \Delta_{q-1}^B=
 \frac{|b_\lambda|^2\nu_0|a_z|^2|1-V_{q-1}|^2}{\kappa_0},
 \qquad
 \Delta_q^B=
 \frac{|b_\lambda|^2\nu_0|g|^2|1-V_q|^2}{s_2\kappa_1},\\
 \boxed{\Phi_{B,q}=
 \frac{2|b_\lambda|^2\nu_0}{\kappa_0}
 \Re\left[-i\overline{a_z}g
 (1-\overline V_{q-1})(1-V_q)\right].}
\end{gathered}
\end{equation}
Every equation follows by inserting the explicitly computed complex
numerators into \eqref{oid:observedloss} and
\eqref{oid:quotientcurrent}. It is a concrete receiving map for the
note's equations (5)--(7); it preserves the still actual
period-dependent responses and restriction determinants.
It does not assign their signs from a full-class energy.
The additional determinant relation \eqref{oid:newallocation} is
especially explicit at this first step:
\begin{equation}\label{oid:lowallocation}
 \boxed{r_q^B=d_0+(1-d_0)\pi_{1,K,q},\qquad
 \kappa_0=\frac{d_0}{d_0+(1-d_0)\pi_{1,K,q}}.}
\end{equation}

% Mathematical section merged verbatim into ORE_INDEPENDENT_DERIVATION.tex.
\section{The actual conductor evaluates the first logarithmic allocation}

The domain of this section is the original simple-packet conductor
domain, $k\equiv1\pmod4$, $k\ge9$, and its fixed actual period
$\varpi$, with the original finite root-comparison guards retained
where the LRS estimates below are used. Its full observation is the
one in inherited FEX1--3~\cite{InheritedFEX}; in particular its
proved conductor identity is
\begin{equation}\label{oid:conductordomain}
 K=\ker\Lambda_k\subset\ker C,
 \qquad C[P]=\bigl((\mathcal T_AP)(\beta'_{\alpha'})\bigr)_{\alpha'},
 \qquad \chi_{k-8}\mid\mathcal T_AP\quad([P]\in K),
\end{equation}
where all lower roots $\beta'_{\alpha'}$ are retained.
The exact receiving map from the original observation is
$\overline C:B\to\C^{q'}$, defined by
$\overline C(\Lambda_kx)=Cx$. It is well-defined because
$\ker\Lambda_k\subset\ker C$, is linear, and is the unique map
with $C=\overline C\Lambda_k$ because $\Lambda_k$ is onto $B$.
Write
\[
 q'=(k-7)^2,\quad c'=c-4,\quad
 \beta_{8,r,s}=4+(2r-8)\delta+i(2s-8)\gamma,
 \quad 0\le r,s\le8,
\]
and retain the complete conductor
\begin{equation}\label{oid:conductoroperator}
 \begin{gathered}
 (\mathcal T_AP)(S')=\sum_{r,s=0}^8a_{rs}P(S'+\beta_{8,r,s}),
 \quad \mu_j=\sum_{r,s=0}^8a_{rs}\beta_{8,r,s}^j,\\
 v_A=\min\{j:\mu_j\ne0\}\le80,
 \qquad \mathcal C_A=\frac{v_A!}{\mu_{v_A}}\mathcal T_A.
 \end{gathered}
\end{equation}
This is the original displayed conductor scalar. It does not change
the arithmetic source or the observation. Taylor expansion of each
shift is finite on a polynomial, so
$\mathcal C_A=\sum_{j\ge v_A}v_A!\mu_j
\partial_S^j/(\mu_{v_A}j!)$. Therefore on a polynomial of degree
$N\ge v_A$ its leading coefficient is exactly $(N)_{v_A}$ times
the original leading coefficient.

The divisibility in \eqref{oid:conductordomain} holds for every
representative of the same class, including the attained one.
Indeed at a lower root the shifted point
$\beta'_{u',w'}+\beta_{8,r,s}=\beta_{k,u'+r,w'+s}$ is an original
upper root. Thus every summand of
$\mathcal T_A(\chi H)(\beta'_{u',w'})$ is zero, for every polynomial
$H$. Simplicity of all lower roots then proves
$\chi_{k-8}\mid\mathcal T_A(\chi H)$. This verifies the exact map
on quotient classes rather than assuming its invariance under the
minimum relation.

Let $\sigma$ be the original order-one Gamma comparison measure,
in the physical coordinate $S=c+iy$ (and $S'=c'+iy$ in the target).
Its complete density and monic norms are
\[
 d\sigma(y)=\frac1{2\pi}|\Gamma(1/4+iy/2)|^2\,dy,\qquad
 \gamma_n=\sqrt{2\pi}\,n!(1/2)_n
 =\sqrt{2\pi}\,\frac{\Gamma(2n+1)}{4^n}.
\]
Retain the original whole-polynomial comparison, valid through degree
$2q$,
\begin{equation}\label{oid:comparison}
 \ell_k\norm P_{\sigma,c}^2\le\norm P_{m_k}^2
 \le u_k\norm P_{\sigma,c}^2,
 \qquad \log(u_k/\ell_k)=O_h(k+\log q).
\end{equation}
This comparison is applied before minimization. It is not a
pointwise lower bound for the arithmetic density.

For clarity, the forward conductor bound of FEX9--11 is reproduced
in the order-one source needed here. In its original monic Gamma
basis $g_n$,
\[
 \sum_{n\ge0}\frac{g_n(y)}{n!}z^n
 =(1+z^2)^{-1/4}e^{y\arctan z},\qquad
 w_n=\frac{n!}{(1/2)_n},\qquad \varphi_n=g_n/\sqrt{\gamma_n}.
\]
The displayed generating identity is also exactly DLMF 18.23.7
\cite{DLMF18} with
$g_n(y)=n!P_n^{(1/4)}(y/2;\pi/2)$: its two factors
$(1-iz)^{-1/4+iy/2}(1+iz)^{-1/4-iy/2}$ multiply to the displayed
right side. This identifies the physical coordinate and monic
coefficient without altering the source norm.
Differentiating the coefficient identity and using
$\arctan z=\sum_{j\ge1,\ j\ {\rm odd}}(-1)^{(j-1)/2}z^j/j$
gives
\[
 \partial_y\varphi_n=
 \sum_{\substack{1\le j\le n\\j\ {\rm odd}}}
 \frac{(-1)^{(j-1)/2}}{j}\sqrt{\frac{w_n}{w_{n-j}}}\,
 \varphi_{n-j}.
\]
The source mass cancels from these coefficient ratios only.
For $n>m\ge1$, the product of
$j(j-1)\le(j-1/2)^2$ gives $w_n/w_m\le\sqrt{n/m}$;
the same product with $w_1=2$ gives $w_n\le2\sqrt n$.
In the absolute derivative matrix use weights $w_n^{-1/2}$.
Its weighted row sums are at most
$\mathsf H_N=\sum_{j=1}^N1/j$. Its weighted column sum at $n$ is
at most $\mathsf H_n+4$, since for $1\le m<n$
\[
 \frac{\sqrt{n/m}-1}{n-m}
 =\frac1{\sqrt m(\sqrt n+\sqrt m)}\le\frac1{\sqrt{nm}},
 \qquad \sum_{m=1}^{n-1}\frac1{\sqrt{nm}}\le2,
\]
and the $m=0$ term is at most $2/\sqrt n\le2$.
Weighted Cauchy--Schwarz on each row, then summation of the column
bound, proves
$\norm{\partial_y}\le\sqrt{\mathsf H_N(\mathsf H_N+4)}
\le\mathsf H_N+2$. The finite exponential of this nilpotent
derivative is translation, so
$\norm{f(y+z)}_\sigma\le e^{|z|(\mathsf H_N+2)}\norm f_\sigma$.
Since the original center displacement is four, triangle inequality
in \eqref{oid:conductoroperator} proves the full retained bound
\begin{equation}\label{oid:forward}
 \begin{gathered}
 \norm{\mathcal C_AP}_{\sigma,c-4}\le M_N\norm P_{\sigma,c},\\
 M_N=\max\left\{1,\frac{v_A!}{|\mu_{v_A}|}
 \sum_{r,s=0}^8|a_{rs}|
 e^{|\beta_{8,r,s}-4|(\mathsf H_N+2)}\right\}.
 \end{gathered}
\end{equation}
No coefficient of the conductor or part of the center displacement
has been suppressed.

Let $\nu^-_{r,1}$ be the original complete monic lower-root minimum:
the minimum of
$\norm{\chi_{k-8}H}_{\sigma,c-4}^2$ over all monic $H$ of degree
$r$. At $N=q-1,q$ put $r_N=N-v_A-q'\ge0$ and let
$P=T_Nx$ for $x\in K$. If its $S^N$ coefficient is $a_N(P)\ne0$,
the polynomial $\mathcal C_AP$ has degree $N-v_A$, is divisible
by the full $\chi_{k-8}$, and has leading coefficient
$(N)_{v_A}a_N(P)$. Therefore
\begin{equation}\label{oid:kernelcoefficient}
 \norm{x}_{G_N}^2=\norm P_{m_k}^2
 \ge\ell_k M_N^{-2}(N)_{v_A}^2
 \nu^-_{r_N,1}|a_N(P)|^2.
\end{equation}
When $a_N(P)=0$ the same inequality is immediate; no nonzero kernel
projection is presumed. The original boundary identity
\eqref{oid:leading} gives
$b_N^*G_Nx=\omega_Na_N(P)$.
For $K=0$ the desired kernel bound is immediate. For $K\ne0$, the
supremum of $|b_N^*G_Nx|^2/\norm x_{G_N}^2$ over nonzero
$x\in K$ is $\norm{P_Nb_N}_{G_N}^2$: Cauchy--Schwarz proves
the upper bound and $x=P_Nb_N$ attains it when this vector is
nonzero, while both sides vanish otherwise.
Also $\omega_N\le u_k\gamma_N$ by applying
\eqref{oid:comparison} to the monic Gamma minimizer.
The exact boundary norms already proved are
\[
 b_N^*G_Nb_N=\omega_N\varepsilon_N^\partial,
 \qquad\varepsilon_{q-1}^\partial=1,\quad
 \varepsilon_q^\partial=1-d_0.
\]
Division only by these positive original factors now proves the
source's conductor inequality, here with its exact receiver:
\begin{equation}\label{oid:finitepi}
 \boxed{\pi_{1,K,N}\le\mathcal B_N:=
 \min\left\{1,\frac{u_k}{\ell_k}
 \frac{M_N^2\gamma_N}
 {\varepsilon_N^\partial(N)_{v_A}^2\nu^-_{r_N,1}}\right\},
 \qquad N=q-1,q.}
\end{equation}
It is a bound on the entire actual kernel of the full original
observation, with all period dependence still present in its
conductor constants.

Combining \eqref{oid:finitepi} at $N=q$ with
\eqref{oid:lowallocation} gives the following new finite allocation
of the first source increment:
\begin{equation}\label{oid:finiteallocation}
 \boxed{\begin{gathered}
 d_0\le r_q^B\le d_0+(1-d_0)\mathcal B_q,\\
 \frac{d_0}{d_0+(1-d_0)\mathcal B_q}\le\kappa_0\le1,\\
 0\le-\log\kappa_0
 \le\log\left(1+\frac{1-d_0}{d_0}\mathcal B_q\right),\\
 -\log\bigl(d_0+(1-d_0)\mathcal B_q\bigr)
 \le-\log r_q^B\le-\log d_0,\\
 -\log\kappa_0-\log r_q^B=-\log d_0.
 \end{gathered}}
\end{equation}
The lower-root minimum in $\mathcal B_q$ can be replaced by its
explicit finite scalar lower enclosure in inherited LRS3--4
\cite{InheritedLRS}. More precisely let
$\underline\tau_{q',r}$ be that original lower endpoint for
$\log\nu^-_{r,1}-\log\gamma_{q'+r}$. Then the same bound holds
with $\mathcal B_N$ replaced by
\begin{equation}\label{oid:evaluatedfinitepi}
 \widehat{\mathcal B}_N=
 \min\left\{1,\frac{u_k}{\ell_k}
 \frac{M_N^2\gamma_N}
 {\varepsilon_N^\partial(N)_{v_A}^2\gamma_{N-v_A}}
 e^{-\underline\tau_{q',r_N}}\right\}.
\end{equation}
This is a proved finite substitution: the lower endpoint implies
$\nu^-_{r_N,1}\ge\gamma_{N-v_A}
e^{\underline\tau_{q',r_N}}$, with its original root guards and
source mass. It introduces no unknown observation determinant.

For the leading rate, inherited LRS3, LRS11 and LRS14
\cite{InheritedLRS} give on those original finite guards
\begin{equation}\label{oid:LRS}
 \begin{gathered}
 \log\nu^-_{r,1}=\log\gamma_{q'+r}+2q'\psi(r/q')+O_h(\log q),\\
 \psi(t)=a_0+\frac t4\log t+(a_0-\tfrac14)t+O(t^{3/2}),
 \qquad a_0=\log(4/\pi).
 \end{gathered}
\end{equation}
These are the inherited complete lower-root norm estimates; they
are used for the actual $r_N=N-v_A-q'$, not a changed lower cutoff.
Here $q-q'=16k-48$, and
$r_{q-1}=16k-49-v_A$, $r_q=16k-48-v_A$. Thus
\[
 \frac{r_N}{q'}=\frac{16}{k}+O(k^{-2}),\qquad
 \frac{r_N}{2}\log\frac{r_N}{q'}=-8k\log k+O(k),\qquad
 \frac{r_N^{3/2}}{\sqrt{q'}}=O(\sqrt k).
\]
Inserting these three quantities into \eqref{oid:LRS} proves
\begin{equation}\label{oid:profile}
 2q'\psi(r_N/q')=2qa_0-8k\log k+O_h(k).
\end{equation}
The remaining monic factor has an exact product, retaining every
factor from the original derivative:
\[
 \frac{\gamma_N}{\gamma_{N-v_A}(N)_{v_A}^2}
 =\prod_{j=0}^{v_A-1}\frac{N-j-1/2}{N-j}\le1.
\]
For fixed original conductor, \eqref{oid:forward} gives
$\log M_N=O_{h,\varpi}(\log q)$. The bound $v_A\le80$
is fixed along the original degree family.

For completeness the first-step determinant rate follows from the
full physical polynomial rather than an observation inference.
Let $R=k\sqrt{\delta^2+\gamma^2}$ and $\alpha_0=\pi/2$.
The original comparison density has fixed constants $c_G,C_G>0$
with
$c_Ge^{-\alpha_0|y|}/\sqrt{1+|y|}\le\sigma(y)
\le C_Ge^{-\alpha_0|y|}/\sqrt{1+|y|}$.
For the actual degree-$q$ monic $Q$, whose full roots lie in
$|y|\le R$, direct integration of
$(|y|-R)^q\le |Q(y)|\le(|y|+R)^q$ on the lower-tail and full
domains respectively gives
\[
 \mathscr L_q(R)\le\norm Q_\sigma^2\le\mathscr U_q(R),
\]
where the lower inequality is integrated only over $|y|\ge R+1$ and
\begin{align*}
 \mathscr L_q(R)&=\frac{2c_Ge^{-\alpha_0R}}{\sqrt{R+2}}
 \alpha_0^{-2q-1/2}\Gamma(2q+\tfrac12,\alpha_0),\\
 \mathscr U_q(R)&=2C_Ge^{\alpha_0R}
 \alpha_0^{-2q-1}\Gamma(2q+1).
\end{align*}
For the lower bound use $s=|y|-R\ge1$ and
$1+R+s\le(R+2)s$; for the upper bound use $s=|y|+R$ and
enlarge the positive integral to $s\ge0$. Monic minimization of
\eqref{oid:comparison} gives
$\ell_k\gamma_q\le\omega_q\le u_k\gamma_q$ and hence
\begin{equation}\label{oid:finited0}
 \frac{\ell_k\gamma_q}{u_k\mathscr U_q(R)}
 \le d_0\le
 \frac{u_k\gamma_q}{\ell_k\mathscr L_q(R)}.
\end{equation}
Stirling's formula and the fixed-shift Gamma-ratio expansion
\cite[\S5.11(i), (iii)]{DLMF5} give
$\log\mathscr U_q(R)-\log\gamma_q
=2q\log(4/\pi)+O_h(k+\log q)$, and the same with
$\mathscr L_q$. The incomplete Gamma tail changes its logarithm
by $o(1)$: the omitted integral from zero to $\alpha_0$ is at most
$\alpha_0^{2q+1/2}/(2q+1/2)$, while Stirling gives the complete
Gamma factor. Therefore
\begin{equation}\label{oid:d0rate}
 \log d_0=-2qa_0+O_h(k+\log q),\qquad
 -\log(1-d_0)=o(1).
\end{equation}
The DLMF expansion is used only for these explicitly displayed
Gamma factors; all physical root and source comparison errors remain.

Equations \eqref{oid:finitepi}, \eqref{oid:comparison},
\eqref{oid:profile} and \eqref{oid:d0rate} now give the source's
actual conductor estimate, with its proof carried into this receiver:
\begin{equation}\label{oid:pirate}
 \pi_{1,K,N}\le\mathcal B_N
 \le\exp\{-2qa_0+8k\log k+O_{h,\varpi}(k)\},
 \qquad N=q-1,q.
\end{equation}
Finally divide the upper bound at $q$ by the lower estimate for
$d_0$ in \eqref{oid:d0rate}, and insert it in
\eqref{oid:finiteallocation}. This proves the new allocation
\begin{equation}\label{oid:firstallocation}
 \boxed{\begin{gathered}
 0\le-\log\frac{\det H_q}{\det H_{q-1}}
 \le8k\log k+O_{h,\varpi}(k),\\
 -\log\frac{\det Q_q}{\det Q_{q-1}}
 =2q\log(4/\pi)+O_{h,\varpi}(k\log k),\\
 \lim_{k\to\infty}\frac{-\log(\det H_q/\det H_{q-1})}{q}=0,
 \qquad
 \lim_{k\to\infty}\frac{-\log(\det Q_q/\det Q_{q-1})}{q}
 =2\log(4/\pi).
 \end{gathered}}
\end{equation}
This evaluates the leading allocation of the single original
increment from $q-1$ to $q$. It neither sums an unproved bound over
the other source increments nor assigns either full two-window
allocation or an individual projected-current sign. The finite
inequality \eqref{oid:finiteallocation}, or its explicit lower-root
enclosure \eqref{oid:evaluatedfinitepi}, remains the receiving formula
before passage to the limit.

\section{Precisely what has been checked and what is additional}

ORE1--5 are the retained definitions, attained observation identity
and source increment proved in \eqref{oid:objects}--\eqref{oid:r}
and \eqref{oid:observation}--\eqref{oid:allocations}.
The mathematical contents of ORE6--10 follow from
\eqref{oid:allocations}--\eqref{oid:response}. The root ORE6 has
a typographical error: the literal string \texttt{,qquad} must be
\texttt{,\textbackslash qquad}. This changes no mathematical factor;
the parent integration subsequently corrected that typography.
ORE11--17 follow from \eqref{oid:products}--\eqref{oid:p1transport},
ORE18--19 from \eqref{oid:fullchart}--\eqref{oid:windows}, and
ORE20--23 from \eqref{oid:observedloss} and
\eqref{oid:quotientcurrent}--\eqref{oid:finitebound}. No correction
of a sign, conjugation, cutoff, or positive denominator is needed in
ORE1--23 on its stated original simple-quartet eigenclass domain.

The supplied low-cutoff note explicitly contains the first and second
energy formulas \eqref{oid:lowfirst} and \eqref{oid:lowsecond}, the
two attained lifts through cutoff $q$, asymptotic estimates for those
energies, a separate conductor bound for the actual kernel overlap,
and the one-step class norm ratio in its equation (17). It does not
contain the following ORE receiving identities:
\begin{enumerate}
\item ORE6--7: the exact restriction/observation determinant allocation
and its two logarithmic increments for the same source step.
\item ORE9--12: the quotient numerator update, consecutive transport
$U_{N+1}=1-(1-V_N)/d_{K,N}$, and the consecutive response products
with their actual earlier conjugation and determinant denominator.
\item ORE14--16: the exact period-dependent observed-energy loss and
the associated update of $\theta_N$. The note's full-class energy
ratio does not by itself determine this quotient energy.
\item ORE17: the general propagation of $p_{2,N}$ into $p_{1,N+1}$.
The note's explicit low-cutoff expressions satisfy its first instance
exactly, but do not state this all-cutoff morphism.
\item ORE18--19: the explicit two-window weighted accumulation of the
same determinant and energy increments, with the full conductor/native
kernel-metric receiver. These identities do not assert new numerical
high-cutoff values.
\item ORE20--23: the undivided observed numerators and consecutive
observed drops computing the current, its phase, its finite variation
bound, and its exact original recurrence/determinant coefficient.
\end{enumerate}
ORE1--3 and ORE8 retain the shared original source, projection and
response definitions. ORE4 is a derived finite minimum identity and
ORE13 its general full-class energy consequence; neither should be
described as an independent evaluation of an observation response.
The note's first low-cutoff energy change is precisely their
specialization \eqref{oid:loweta} at $q-1$.

Additional to both the root ORE identities and the note's equations
(5)--(7), this derivation proves the minimum-lift/projector update
\eqref{oid:projector}, the next-overlap determinant morphism
\eqref{oid:newallocation}, the strict original bound in
\eqref{oid:Eloss}, the exact next source minimum and norm in
\eqref{oid:loweta}, the low-cutoff full-current transport
\eqref{oid:lowcurrent}--\eqref{oid:lowmean}, and the explicit
observed-current and determinant receivers
\eqref{oid:lowobservation}--\eqref{oid:lowallocation}, together with
the finite and asymptotic single-increment allocations
\eqref{oid:finiteallocation} and \eqref{oid:firstallocation}.
The latter use the complete conductor bound proved above and the
explicitly cited inherited LRS lower-root estimates. They are a new
consequence for the actual observation, not merely the supplied
note's full-class energy calculation. The note's absolute-action
and marked-family calculations are outside this derivation.

\begin{thebibliography}{9}
\bibitem{Hager} W.~W. Hager, ``Updating the Inverse of a Matrix'',
\emph{SIAM Review} 31 (1989), 221--239,
\href{https://doi.org/10.1137/1031049}{doi:10.1137/1031049}.
\bibitem{DLMF5} R.~A. Askey and R. Roy, ``Gamma Function'',
Chapter 5 in the NIST Digital Library of Mathematical Functions,
\S5.11, ``Asymptotic Expansions'', especially (i), (ii), and (iii),
\url{https://dlmf.nist.gov/5.11}.
\bibitem{DLMF18} T.~H. Koornwinder, R. Wong, R. Koekoek,
R.~F. Swarttouw and W.~P. Reinhardt, ``Orthogonal Polynomials'',
Chapter 18 in the NIST Digital Library of Mathematical Functions,
\S18.23, equation (18.23.7),
\url{https://dlmf.nist.gov/18.23#E7}. The DLMF source note attributes
this generating identity to Ismail (2009), equation (5.9.3).
\bibitem{InheritedFEX} Split-Zero programme, original complete
fixed-period conductor proof, FEX1--11, in
\texttt{14\_CURRENT\_KERNEL\_AND\_MULTIPLICITY\_PROOFS.tex}.
The observation, source, conductor coefficients, and Gamma derivative
identity in this derivation are those retained in that source.
\bibitem{InheritedLRS} Split-Zero programme, ``Literal covariance,
logarithmic endpoint term, and the complete action'', LRS3--4 and
LRS11--14, in \texttt{LRC\_SCALAR\_RECONSTRUCTION.tex}.
The original finite root guards and complete lower-root norm estimates
are retained in their application above.
\end{thebibliography}

\end{document}
